\section{Noetherian local rings}
\label{section-Noetherian-local}

\noindent
In all of this section $(R, \mathfrak m, \kappa)$ is a Noetherian local ring.
We develop some theory on Hilbert functions of modules in this section.
Let $M$ be a finite $R$-module. We define the {\it Hilbert function}
of $M$ to be the function
$$
\varphi_M : n
\longmapsto
\text{length}_R(\mathfrak m^nM/{\mathfrak m}^{n + 1}M)
$$
defined for all integers $n \geq 0$. Another important invariant is the
function
$$
\chi_M : n
\longmapsto
\text{length}_R(M/{\mathfrak m}^{n + 1}M)
$$
defined for all integers $n \geq 0$.
Note that we have by Lemma \ref{lemma-length-additive}
that
$$
\chi_M(n) = \sum\nolimits_{i = 0}^n \varphi_M(i).
$$
There is a variant of this construction which uses an ideal of definition.

\begin{definition}
\label{definition-ideal-definition}
Let $(R, \mathfrak m)$ be a local Noetherian ring.
An ideal $I \subset R$ such that $\sqrt{I} = \mathfrak m$ is called
{\it an ideal of definition of $R$}.
\end{definition}

\noindent
Let $I \subset R$ be an ideal of definition.
Because $R$ is Noetherian this means that
$\mathfrak m^r \subset I$ for some $r$, see Lemma
\ref{lemma-Noetherian-power}. Hence any finite $R$-module
annihilated by a power of $I$ has a finite length, see Lemma
\ref{lemma-length-finite}.
Thus it makes sense to define
$$
\varphi_{I, M}(n) = \text{length}_R(I^nM/I^{n + 1}M)
\quad\text{and}\quad
\chi_{I, M}(n) = \text{length}_R(M/I^{n + 1}M)
$$
for all $n \geq 0$. Again we have that
$$
\chi_{I, M}(n) = \sum\nolimits_{i = 0}^n \varphi_{I, M}(i).
$$

\begin{lemma}
\label{lemma-differ-finite}
Suppose that $M' \subset M$ are finite $R$-modules
with finite length quotient. Then there exists a
constants $c_1, c_2$ such that for all $n \geq c_2$ we have
$$
c_1 + \chi_{I, M'}(n - c_2) \leq \chi_{I, M}(n) \leq
c_1 + \chi_{I, M'}(n)
$$
\end{lemma}

\begin{proof}
Since $M/M'$ has finite length there is a $c_2 \geq 0$ such that
$I^{c_2}M \subset M'$. Let $c_1 = \text{length}_R(M/M')$.
For $n \geq c_2$ we have
\begin{eqnarray*}
\chi_{I, M}(n)
& = &
\text{length}_R(M/I^{n + 1}M) \\
& = &
c_1 + \text{length}_R(M'/I^{n + 1}M) \\
& \leq &
c_1 + \text{length}_R(M'/I^{n + 1}M') \\
& = &
c_1 + \chi_{I, M'}(n)
\end{eqnarray*}
On the other hand, since $I^{c_2}M \subset M'$,
we have $I^nM \subset I^{n - c_2}M'$ for $n \geq c_2$.
Thus for $n \geq c_2$ we get
\begin{eqnarray*}
\chi_{I, M}(n)
& = &
\text{length}_R(M/I^{n + 1}M) \\
& = &
c_1 + \text{length}_R(M'/I^{n + 1}M) \\
& \geq &
c_1 + \text{length}_R(M'/I^{n + 1 - c_2}M') \\
& = &
c_1 + \chi_{I, M'}(n - c_2)
\end{eqnarray*}
which finishes the proof.
\end{proof}

\begin{lemma}
\label{lemma-hilbert-ses}
Suppose that $0 \to M' \to M \to M'' \to 0$
is a short exact sequence of finite $R$-modules.
Then there exists a submodule $N \subset M'$ with
finite colength $l$ and $c \geq 0$ such that
$$
\chi_{I, M}(n) = \chi_{I, M''}(n) + \chi_{I, N}(n - c) + l
$$
and
$$
\varphi_{I, M}(n) = \varphi_{I, M''}(n) + \varphi_{I, N}(n - c)
$$
for all $n \geq c$.
\end{lemma}

\begin{proof}
Note that $M/I^nM \to M''/I^nM''$ is surjective
with kernel $M' / M' \cap I^nM$. By the Artin-Rees
Lemma \ref{lemma-Artin-Rees} there exists a
constant $c$ such that $M' \cap I^nM =
I^{n - c}(M' \cap I^cM)$. Denote $N = M' \cap I^cM$.
Note that $I^c M' \subset N \subset M'$.
Hence $\text{length}_R(M' / M' \cap I^nM)
= \text{length}_R(M'/N) + \text{length}_R(N/I^{n - c}N)$ for $n \geq c$.
From the short exact sequence
$$
0 \to M' / M' \cap I^nM \to M/I^nM \to M''/I^nM'' \to 0
$$
and additivity of lengths (Lemma \ref{lemma-length-additive})
we obtain the equality
$$
\chi_{I, M}(n - 1)
=
\chi_{I, M''}(n - 1)
+
\chi_{I, N}(n - c - 1)
+
\text{length}_R(M'/N)
$$
for $n \geq c$. We have
$\varphi_{I, M}(n) = \chi_{I, M}(n) - \chi_{I, M}(n - 1)$
and similarly for the modules $M''$ and $N$. Hence
we get $\varphi_{I, M}(n) = \varphi_{I, M''}(n) + \varphi_{I, N}(n-c)$ for
$n \geq c$.
\end{proof}

\begin{lemma}
\label{lemma-hilbert-change-I}
Suppose that $I$, $I'$ are two ideals of definition
for the Noetherian local ring $R$. Let $M$ be a
finite $R$-module. There exists a constant $a$ such that
$\chi_{I, M}(n) \leq \chi_{I', M}(an)$ for $n \geq 1$.
\end{lemma}

\begin{proof}
There exists an integer $c \geq 1$ such that $(I')^c \subset I$.
Hence we get a surjection $M/(I')^{c(n + 1)}M \to M/I^{n + 1}M$.
Whence the result with $a = 2c - 1$.
\end{proof}

\begin{proposition}
\label{proposition-hilbert-function-polynomial}
Let $R$ be a Noetherian local ring. Let $M$ be a finite $R$-module.
Let $I \subset R$ be an ideal of definition.
The Hilbert function $\varphi_{I, M}$ and the function
$\chi_{I, M}$ are numerical polynomials.
\end{proposition}

\begin{proof}
Consider the graded ring $S = R/I \oplus I/I^2 \oplus I^2/I^3 \oplus
\ldots = \bigoplus_{d \geq 0} I^d/I^{d + 1}$. Consider the graded
$S$-module $N = M/IM \oplus IM/I^2M \oplus \ldots =
\bigoplus_{d \geq 0} I^dM/I^{d + 1}M$. This pair $(S, N)$ satisfies
the hypotheses of Proposition \ref{proposition-graded-hilbert-polynomial}.
Hence the result for $\varphi_{I, M}$ follows from that proposition and
Lemma \ref{lemma-length-K0}. The result for $\chi_{I, M}$ follows
from this and Lemma \ref{lemma-numerical-polynomial}.
\end{proof}

\begin{definition}
\label{definition-hilbert-polynomial}
Let $R$ be a Noetherian local ring. Let $M$ be a finite $R$-module.
The {\it Hilbert polynomial} of $M$ over $R$ is the element
$P(t) \in \mathbf{Q}[t]$ such that $P(n) = \varphi_M(n)$ for $n \gg 0$.
\end{definition}

\noindent
By Proposition \ref{proposition-hilbert-function-polynomial}
we see that the Hilbert polynomial exists.

\begin{lemma}
\label{lemma-d-independent}
Let $R$ be a Noetherian local ring. Let $M$ be a finite $R$-module.
\begin{enumerate}
\item The degree of the numerical polynomial $\varphi_{I, M}$ is independent
of the ideal of definition $I$.
\item The degree of the numerical polynomial $\chi_{I, M}$ is independent
of the ideal of definition $I$.
\end{enumerate}
\end{lemma}

\begin{proof}
Part (2) follows immediately from Lemma \ref{lemma-hilbert-change-I}.
Part (1) follows from (2) because
$\varphi_{I, M}(n) = \chi_{I, M}(n) - \chi_{I, M}(n - 1)$
for $n \geq 1$.
\end{proof}

\begin{definition}
\label{definition-d}
Let $R$ be a local Noetherian ring and $M$ a finite $R$-module.
We denote {\it $d(M)$} the element of $\{-\infty, 0, 1, 2, \ldots \}$
defined as follows:
\begin{enumerate}
\item If $M = 0$ we set $d(M) = -\infty$,
\item if $M \not = 0$ then $d(M)$ is the degree of the numerical
polynomial $\chi_M$.
\end{enumerate}
\end{definition}

\noindent
If $\mathfrak m^nM \not = 0$ for all $n$, then we see that
$d(M)$ is the degree $+1$ of the Hilbert polynomial of $M$.

\begin{lemma}
\label{lemma-differ-finite-chi}
Let $R$ be a Noetherian local ring. Let $I \subset R$ be an ideal
of definition. Let $M$ be a finite $R$-module
which does not have finite length. If $M' \subset M$ is a submodule
with finite colength, then $\chi_{I, M} - \chi_{I, M'}$
is a polynomial of degree $<$ degree of either polynomial.
\end{lemma}

\begin{proof}
Follows from Lemma \ref{lemma-differ-finite} by elementary calculus.
\end{proof}

\begin{lemma}
\label{lemma-hilbert-ses-chi}
Let $R$ be a Noetherian local ring. Let $I \subset R$ be an ideal of
definition. Let $0 \to M' \to M \to M'' \to 0$ be a short exact sequence
of finite $R$-modules. Then
\begin{enumerate}
\item if $M'$ does not have finite length, then
$\chi_{I, M} - \chi_{I, M''} - \chi_{I, M'}$
is a numerical polynomial of degree $<$ the degree of
$\chi_{I, M'}$,
\item $\max\{ \deg(\chi_{I, M'}), \deg(\chi_{I, M''}) \} = \deg(\chi_{I, M})$,
and
\item $\max\{d(M'), d(M'')\} = d(M)$,
\end{enumerate}
\end{lemma}

\begin{proof}
We first prove (1). Let $N \subset M'$ be as in Lemma \ref{lemma-hilbert-ses}.
By Lemma \ref{lemma-differ-finite-chi} the numerical polynomial
$\chi_{I, M'} - \chi_{I, N}$ has degree $<$ the common degree of
$\chi_{I, M'}$ and $\chi_{I, N}$. By Lemma \ref{lemma-hilbert-ses}
the difference
$$
\chi_{I, M}(n) - \chi_{I, M''}(n) - \chi_{I, N}(n - c)
$$
is constant for $n \gg 0$. By elementary calculus the difference
$\chi_{I, N}(n) - \chi_{I, N}(n - c)$ has degree $<$ the degree of
$\chi_{I, N}$ which is bigger than zero (see above). Putting everything
together we obtain (1).

\medskip\noindent
Note that the leading coefficients of $\chi_{I, M'}$ and $\chi_{I, M''}$ are
nonnegative. Thus the degree of $\chi_{I, M'} + \chi_{I, M''}$ is equal
to the maximum of the degrees. Thus if $M'$ does not have finite
length, then (2) follows from (1). If $M'$ does have finite length, then
$I^nM \to I^nM''$ is an isomorphism for all $n \gg 0$ by Artin-Rees
(Lemma \ref{lemma-Artin-Rees}). Thus $M/I^nM \to M''/I^nM''$ is a
surjection with kernel $M'$ for $n \gg 0$ and we see that
$\chi_{I, M}(n) - \chi_{I, M''}(n) = \text{length}(M')$
for all $n \gg 0$. Thus (2) holds in this case also.

\medskip\noindent
Proof of (3). This follows from (2) except if one of $M$, $M'$, or $M''$
is zero. We omit the proof in these special cases.
\end{proof}



































